% Section 3 -- Compiled Stateflow IR
%
% Goal: introduce the typed IR.  Keep it self-contained -- a reader who
% skipped the background should still be able to see what a plan is,
% what it costs, and how we reject bad plans.

\subsection{Design Goals}\label{sec:ir:goals}

Out-of-core KGE training is a scheduling problem whose decisions --
which partition to evict, which to admit, from where to source the
admit -- happen at runtime but are determined by choices made long
before any GPU kernel runs. We designed Compiled Stateflow around four
goals that fall out of that observation.

\paragraph{G1: Every byte is typed.} Host admits, device-to-device peer
copies, and kept-alive resident pages differ by two or three orders of
magnitude in latency and by an order of magnitude in bandwidth. A
scheduler that cannot distinguish them cannot reason about them. Every
transfer in Stateflow carries a mode tag and a byte count.

\paragraph{G2: Plans are values, not side effects.} A Stateflow plan is
a pure data structure. It is produced by the compiler, handed to a
validator, and only then consumed by the runtime. This separation lets
us enumerate candidates, cost them, apply invariants, and diff plans
across configurations without running the trainer.

\paragraph{G3: Cost is explicit and decomposable.} The scheduler's
objective function is not a scalar black box. A plan's cost decomposes
into admit cost, bucket imbalance, boundary crossings, and lane
imbalance, each attributable to a concrete portion of the plan. This
lets the compiler explain \emph{why} one plan beat another, which in
turn lets us write targeted transformations.

\paragraph{G4: The IR survives Phase changes.} We stage the system in
phases -- single-GPU scheduling, multi-GPU compilation, peer runtime,
bandwidth-aware staggering -- and the IR is the one artifact that
remains stable across them. Phase 3 adds descriptors the runtime does
not yet execute; the final Phase 4b runtime now executes them without
changing the compiler. Phase 5's staggering will introduce new cost
terms without changing the type hierarchy. A type that absorbs
additions without breaking consumers is the whole reason to build an
IR at all.

\subsection{Type Hierarchy}\label{sec:ir:types}

A plan is a nested structure (Figure~\ref{fig:ir:types}). At the top,
\texttt{StateflowPlan} holds one \texttt{LanePlan} per active GPU plus
a top-level list of cross-lane handoffs. Each \texttt{LanePlan} is an
ordered sequence of \texttt{Microstate}s -- the per-lane steps between
round boundaries -- connected by intra-lane \texttt{Handoff}s. A
\texttt{Microstate} names the partitions resident on the device at that
step, the fragments they occupy, and the resident objects (embeddings,
optimizer state) that back them.

\begin{figure}[t]
\centering
\begin{lstlisting}[basicstyle=\ttfamily\footnotesize]
StateflowPlan
  family        : {CUSTOM, HYBRID_COVER, MULTI_GPU}
  variant       : {CANONICAL, GPU_AWARE, REVERSED, ...
                   DISJOINT_ROUNDS, LANE_MATCHED}
  lanes         : [LanePlan]
  cross_lane_handoffs : [HandoffPlan]    # top-level only
  cost          : PlanCostBreakdown
  total_peer_handoff_bytes : int64

LanePlan
  lane_id       : int
  microstates   : [MicrostatePlan]       # length N
  handoffs      : [HandoffPlan]          # length N-1, intra-lane

MicrostatePlan
  round_idx     : int
  resident_partitions : [int]
  fragments     : [FragmentPlan]
  resident_objects    : [ResidentObjectPlan]

HandoffPlan
  mode          : {FULL_RELOAD, ROTATING_OVERWRITE,
                   DELAYED_KEEP_ALIVE, PEER_RELAY}
  partition_id  : int
  src_lane_id   : int     # -1 if host-sourced
  dst_lane_id   : int
  bytes         : int64
  peer_bytes    : int64   # nonzero iff PEER_RELAY

PeerHandoffDescriptor                    # runtime target
  src_lane, dst_lane, src_slot, dst_slot : int
  round_idx, partition_id                : int
  bytes                                  : int64
\end{lstlisting}
\caption{Core Stateflow types. The IR separates intra-lane handoffs
(attached to a \texttt{LanePlan}) from cross-lane handoffs (attached to
the top-level plan). The runtime consumes \texttt{PeerHandoffDescriptor}s
projected from the plan.}
\label{fig:ir:types}
\end{figure}

Two design choices in Figure~\ref{fig:ir:types} are worth calling out.

First, \textbf{cross-lane handoffs live at the top level}, not inside
\texttt{LanePlan.handoffs}. A \texttt{LanePlan} with $N$ microstates
must have exactly $N-1$ intra-lane handoffs; this invariant is checked
by the validator and lets downstream code zip lanes with their own
transitions. Mixing cross-lane handoffs into that list would break the
invariant or require per-entry tagging at every call site. We instead
expose them through \texttt{StateflowPlan.cross\_lane\_handoffs}, a
flat list the runtime indexes by \texttt{round\_idx}.

Second, \texttt{HandoffPlan} carries both \texttt{bytes} and
\texttt{peer\_bytes}. The former is the host-fallback cost (what a
host admit would transfer); the latter is nonzero only for
\texttt{PEER\_RELAY} and captures the device-to-device byte count the
runtime would issue via \texttt{cudaMemcpyPeerAsync}. Keeping them
separate lets the default cost model be host-denominated (safe) while
still reporting peer-executable volume for Phase 4 headroom analysis.

\subsection{Handoff Modes}\label{sec:ir:modes}

Partition residency changes at round boundaries. We identified four
distinct ways a partition can become resident on lane $\ell$ at round
$r{+}1$; each is a separate handoff mode because each has a different
cost and different runtime mechanism.

\begin{description}
\item[\texttt{FULL\_RELOAD}] The partition is fetched from host memory
as if this were a cold admit. Cost: full host-bandwidth transfer of
embeddings and optimizer state. Occurs when no lane held the partition
in round $r$, or when the holder did not share an NVLink/PCIe peer
path to $\ell$.
\item[\texttt{ROTATING\_OVERWRITE}] The partition was resident on $\ell$
in round $r$ and its slot is reused in round $r{+}1$. No transfer; the
cost is zero. Named for the rotating-window layouts that naturally
produce this pattern.
\item[\texttt{DELAYED\_KEEP\_ALIVE}] The partition was resident on $\ell$
in round $r$, is not needed in round $r{+}1$, but is scheduled for
round $r{+}2$. The runtime keeps it pinned to avoid reloading it one
round later. Cost: device memory occupancy for one extra round; no
transfer.
\item[\texttt{PEER\_RELAY}] The partition is resident on lane $\ell'$
in round $r$, not on $\ell$, but needed on $\ell$ in round $r{+}1$. A
device-to-device copy from $\ell'$ to $\ell$ delivers it. Cost:
\texttt{peer\_bytes} at peer bandwidth (substantially below host
bandwidth on modern interconnects). Requires
\texttt{cudaDeviceCanAccessPeer} to be true.
\end{description}

\texttt{PEER\_RELAY} is the mode that makes the multi-GPU case
interesting. The compiler emits it whenever lane matching produces a
cross-lane opportunity; the runtime executes it as
\texttt{cudaMemcpyPeerAsync}. The same descriptor supported both
bring-up regimes: Phase 3 logged it while still executing a host
reload, and the final Phase 4b runtime consumes it directly without
changing the compiler or validator.

\subsection{Cost Model}\label{sec:ir:cost}

Candidate plans are ranked by a weighted sum of four terms:
\begin{align}
\mathrm{cost}(P) =\ & \alpha \cdot \mathrm{admit}(P)
                    + \beta  \cdot \mathrm{bucket}(P) \nonumber \\
                  & + \gamma \cdot \mathrm{boundary}(P)
                    + \delta \cdot \mathrm{imbalance}(P).
\label{eq:cost}
\end{align}

\paragraph{Admit cost.} $\mathrm{admit}(P) = \sum_h b_h$, summed over
handoffs $h$ that cause a host fetch. Measured in bytes; this is the
host-bandwidth-denominated cost of populating microstates whose
partitions were not resident in the preceding round. In default
configuration, \texttt{PEER\_RELAY} is counted as a \emph{host} admit
(its bytes would be loaded from host if the peer path were
unavailable). This is a deliberately conservative choice: until the
selector is made peer-aware, the compiler must not give itself credit
for savings it is not yet optimizing for.

\paragraph{Bucket cost.} $\mathrm{bucket}(P)$ penalizes uneven edge
bucket sizes across microstates. Small buckets underutilize the GPU;
enormous buckets over-fill a single superstep. We weight this term
with a small $\beta \!\approx\! 10^{-7}$ so it acts as a tie-breaker
among plans with similar admit cost.

\paragraph{Boundary cost.} $\mathrm{boundary}(P)$ counts superstate
boundaries -- transitions that force global synchronization across
lanes. In the multi-GPU setting these are NCCL all-reduce points;
each one is a latency cost plus a bandwidth tax on embedding
gradients.

\paragraph{Imbalance cost.} $\mathrm{imbalance}(P) = \max_\ell W_\ell
- \min_\ell W_\ell$, where $W_\ell$ is lane $\ell$'s per-round work
(typically summed bucket bytes). A plan where one lane finishes
substantially before another wastes the gap; this term pushes the
compiler toward lane assignments where the work lands evenly.

\paragraph{Reported metrics, not optimized.} Two quantities are
recorded on every \texttt{StateflowPlan} but not summed into
Equation~\ref{eq:cost}: \texttt{total\_peer\_handoff\_bytes} and
\texttt{total\_cross\_lane\_handoffs}. These characterize the plan's
Phase 4 opportunity (how much host I/O a peer runtime could displace)
without letting the selector claim credit for it during compilation.
Even now that the runtime executes peer copies, the default selector
keeps these as reported metrics rather than optimized terms. A future
peer-aware cost model can sum them into Equation~\ref{eq:cost}
explicitly -- a clean extension point that does not disturb Phases 1--4.

\subsection{Validator Invariants}\label{sec:ir:validator}

A plan is only useful if it is executable. The validator rejects plans
that violate any of the following invariants; no plan leaves the
compiler without passing.

\begin{description}
\item[V1: Structural consistency.] Each lane has
$|\text{handoffs}| = |\text{microstates}| - 1$. Cross-lane handoffs
appear only in \texttt{StateflowPlan.cross\_lane\_handoffs}, never in
\texttt{LanePlan.handoffs}.
\item[V2: Coverage.] For every microstate boundary $(r, r{+}1)$ on
every lane $\ell$, the resident set at $r{+}1$ equals
$(\text{kept}_{\ell}) \cup (\text{host-admitted}) \cup
 (\text{peer-received})$. No partition appears on a lane without a
sourcing handoff.
\item[V3: Peer admit sourcing.] A partition delivered by
\texttt{PEER\_RELAY} must be resident on the source lane in round $r$
and must not be resident on the destination lane in round $r$. Both
ends are checked.
\item[V4: Superstate disjointness.] Within a round, partitions
assigned to the same superstate must be disjoint across lanes.
Violating this would double-count gradient contributions during the
all-reduce.
\item[V5: Handoff consistency.] Every cross-lane handoff with mode
\texttt{PEER\_RELAY} has matching source and destination lanes in the
microstate's resident sets. The byte count agrees with the embedding
layout (Section~\ref{sec:ir:layout}).
\end{description}

Invariants V1--V5 are machine-checked and cheap -- the validator
passes in microseconds on 16-partition plans. They serve two
purposes: they catch compiler bugs at their source (a corrupt plan
cannot reach the runtime), and they document the contract between the
IR and every downstream consumer.

\subsection{Embedding Layout}\label{sec:ir:layout}

Byte counts are not inferable from partition IDs alone: they depend on
embedding dimension, dtype, and whether the trainer uses optimizer
state that doubles the resident footprint. \texttt{PlanEmbeddingLayout}
captures this:

\begin{lstlisting}[basicstyle=\ttfamily\footnotesize]
PlanEmbeddingLayout
  embedding_dim           : int    # e.g. 400
  dtype_size              : int    # bytes per element
  optimizer_state_multiplier : int # 1 or 2
\end{lstlisting}

\noindent
A partition of $n$ rows contributes
$n \cdot \mathtt{embedding\_dim} \cdot \mathtt{dtype\_size}
     \cdot \mathtt{optimizer\_state\_multiplier}$ bytes to any handoff
that transfers it. The layout is populated once at plan construction
from the graph storage metadata, then embedded in every plan so
downstream tooling (cost reporters, visualizers) can produce concrete
byte figures without re-querying the model.

\subsection{A Concrete Example}\label{sec:ir:example}

Consider Twitter (41.6M nodes, $d\!=\!400$, fp32, Adagrad), 16
partitions, buffer capacity 4, two GPUs. Two compiled plans:

\begin{table}[h]
\centering
\small
\begin{tabular}{lrr}
\toprule
 & \textsc{Disjoint\_Rounds} & \textsc{Lane\_Matched} \\
\midrule
selection cost        & $1.499 \times 10^{11}$ & $1.291 \times 10^{11}$ \\
host admits           & 72                     & 62                     \\
cross-lane handoffs   & 8                      & 18                     \\
\texttt{PEER\_RELAY} bytes & 16.66 GB          & 37.49 GB               \\
\texttt{FULL\_RELOAD} count & 10               & 0                      \\
\texttt{ROTATING\_OVERWRITE} & 8               & 18                     \\
boundaries            & 18                     & 18                     \\
\bottomrule
\end{tabular}
\caption{Two compiled Stateflow plans on Twitter 16-partition, 2-GPU.
\textsc{Lane\_Matched} lowers host-denominated cost by 14\% and
eliminates all \texttt{FULL\_RELOAD} events and exposes 37.49 GB of
peer-relay opportunity. Section~\ref{sec:eval} shows the runtime
executing the embeddings half of that traffic with zero fallback.}
\label{tab:ir:example}
\end{table}

Both plans satisfy V1--V5 and are produced by the same compiler under
different variant flags. \textsc{Lane\_Matched} is the output of the
compiler described in Section~\ref{sec:compiler}; \textsc{Disjoint\_Rounds}
is a baseline that ignores cross-lane opportunities. The entire
difference is visible in the IR -- no runtime is needed to see that
\textsc{Lane\_Matched} has eliminated the full reloads. The runtime
results of Section~\ref{sec:eval} then show that these IR-level
differences survive projection and swap-time execution exactly.
